
\begin{frame}{L'optimiseur}

L'optimiseur ORACLE suit une approche classique\,:
\begin{redbullet}
   \item Génération de plusieurs plans d'exécution.
   \item Estimation du coût de chaque plan généré.
   \item Choix du meilleur et exécution.
\end{redbullet}

Tout ceci est automatique, \alert{mais il est possible d'influer,
voire de forcer le plan d'exécution}.
\end{frame}

\begin{frame}{Les chemins d'accès}

Ceux que nous avons déjà vus.


\begin{redbullet}
  \item \alert{Parcours séquentiel} ($FULL\ TABLE\ SCAN$).
  \item \alert{Par adresse} ($ACCESS\ BY\ ROWID$).
 \item \alert{Parcours d'index} ($INDEX\ SCAN$).
\end{redbullet}

\end{frame}

\begin{frame}{Opérations physiques}


Voici les principales\,:
\begin{redbullet}
\item $INTERSECTION$ : intersection de deux ensembles de n-uplet.
\item $CONCATENATION$ : union de deux ensembles.
\item $FILTER$ : élimination de n-uplets (sélection).
\item $PROJECTION$ : opération de l'algèbre relationnelle.
\end{redbullet}

D'autres opérations sont liées aux algorithmes de jointures.

\end{frame}

\begin{frame}{Algorithmes de jointure sous ORACLE}


ORACLE utilise trois algorithmes de jointures\,:
\begin{redbullet}
  \item \alert{Boucles imbriquées} quand il y a au moins un index.\\
	Opération \texttt{NESTED\ LOOP}.
   \item  \alert{Tri/fusion} quand il n'y a pas d'index.\\
	Opération \texttt{SORT} et \texttt{MERGE}.

   \item  \alert{Jointure par hachage} quand il n'y a pas d'index.\\
	Opération  \texttt{HASH JOIN}

\end{redbullet}

\end{frame}


\begin{frame}{L'outil EXPLAIN}

L'outil EXPLAIN donne le plan d'exécution d'une requ\^ete. La
description comprend\,:
\begin{redbullet}
  \item Le chemin d'accès utilisé.
  \item Les opérations physiques (tri, fusion, intersection, ...).
  \item L'ordre des opérations. 
\end{redbullet}

Il est représentable par un arbre.

\end{frame}


\subsection{Optimisation dans ORACLE\,: exemples}

\begin{frame}{Rappel du schéma}
\begin{redbullet}
  \item Film (\textbf{idFilm}, titre, année, genre, résumé,
                                \textit{idMES}, \textit{codePays})
 \item Artiste (\textbf{idArtiste}, nom, prénom, annéeNaissance)
  \item Role (\textbf{\textit{idActeur}, \textit{idFilm}}, nomRôle)
   \item Internaute (\textbf{email}, nom, prénom, région)
  \item Notation (\textbf{\textit{email}, \textit{idFilm}}, note)
   \item Pays (\textbf{code}, nom, langue)
\end{redbullet}

\end{frame}

\begin{frame}[fragile]{Sélection sans index}

La requête\,:

\begin{minted}[frame=leftline,
               baselinestretch=.9,
               fontsize=\footnotesize,
               framesep=2mm]{sql}
     explain plan 
     set statement_id='SelSansInd' for
     select *
     from Film
     where titre = 'Vertigo' 
\end{minted}

\vfill

Le résultat de EXPLAIN\,:
\begin{minted}[frame=leftline,
               baselinestretch=.9,
               fontsize=\footnotesize,
               framesep=2mm]{sql}
    0 SELECT STATEMENT
      1 TABLE ACCESS FULL FILM   
\end{minted}

\end{frame}

\begin{frame}{Plan d'exécution}

On ne fait pas plus simple\,!

\figSlideWithSize{planOra-selsansind}{3cm}

\end{frame}

\begin{frame}[fragile]{Sélection avec index}

La requête\,:

\begin{minted}[frame=leftline,
               baselinestretch=.9,
               fontsize=\footnotesize,
               framesep=2mm]{sql}
     explain plan 
     set statement_id='SelInd' for
     select *
     from Film
     where id=21;
\end{minted}

\vfill

Le résultat de EXPLAIN\,:
\begin{minted}[frame=leftline,
               baselinestretch=.9,
               fontsize=\footnotesize,
               framesep=2mm]{sql}
    0 SELECT STATEMENT
     1 TABLE ACCESS BY ROWID FILM
      2 INDEX UNIQUE SCAN IDX-FILM-ID 
\end{minted}

\end{frame}


%%%%%%%%%%%%%%%%%%%%%%%%%%%%%%%%%%%%%%%%%%%%%%%%%%%%%%%
\begin{frame}{Plan d'exécution}

Accès à l'index, puis à la table.

\figSlideWithSize{planOra-selind}{6cm}

\end{frame}


\begin{frame}[fragile]{Jointure avec index}

La requête\,:

\begin{minted}[frame=leftline,
               baselinestretch=.9,
               fontsize=\footnotesize,
               framesep=2mm]{sql}
    explain plan 
    set statement_id='JoinIndex' for
    select titre, nom, prenom
    from   Film f, Artiste a
    where f.id_realisateur = a.id;
\end{minted}

\vfill

Le résultat de EXPLAIN\,:
\begin{minted}[frame=leftline,
               baselinestretch=.9,
               fontsize=\footnotesize,
               framesep=2mm]{sql}
    0 SELECT STATEMENT
      1 NESTED LOOPS
        2 TABLE ACCESS FULL FILM
        3 TABLE ACCESS BY ROWID ARTISTE
          4 INDEX UNIQUE SCAN IDXARTISTE
\end{minted}

\end{frame}


\begin{frame}{Plan d'exécution (déjà vu)}

\figSlideWithSize{planEx-indexedjoin}{6cm}
\end{frame}

\begin{frame}[fragile]{Jointure avec index et sélection}

La requête (avec index sur le nom des artistes)\,:

\begin{minted}[frame=leftline,
               baselinestretch=.9,
               fontsize=\footnotesize,
               framesep=2mm]{sql}
    explain plan
    set statement_id='JoinSelIndex' for
    select nom_role
    from   Role r, Artiste a
    where  r.id_acteur = a.id
    and    nom = 'Pacino'; 
\end{minted}

\vfill

Le résultat de EXPLAIN\,:
\begin{minted}[frame=leftline,
               baselinestretch=.9,
               fontsize=\footnotesize,
               framesep=2mm]{sql}
    0 SELECT STATEMENT
      1 NESTED LOOPS
        2 TABLE ACCESS BY ROWID ARTISTE
          3 INDEX RANGE SCAN IDX-NOM
        4 TABLE ACCESS BY ROWID ROLE
          5 INDEX RANGE SCAN IDX-ROLE
\end{minted}

\end{frame}

\begin{frame}{Plan d'exécution}

\figSlideWithSize{planOra-joinselindex}{10cm}

\end{frame}

%%%%%%%%%%%%%%%%%%%%%%%%%%%%%%%%%%%%%%%%%%%%%%%%%%%%%%%
\begin{frame}[fragile]{Jointure sans index}

Requête: 
\begin{minted}[frame=leftline,
               baselinestretch=.9,
               fontsize=\footnotesize,
               framesep=2mm]{sql}
    explain plan set 
    statement_id='JoinSansIndex' for
    select nom, prenom
    from Film f, Artiste a
    where f.annee  = a.annee_naissance
    and   titre = 'Vertigo';
\end{minted}

\vfill

Le résultat de EXPLAIN\,:
\begin{minted}[frame=leftline,
               baselinestretch=.9,
               fontsize=\footnotesize,
               framesep=2mm]{sql}
   0 SELECT STATEMENT
     1 MERGE JOIN
       2 SORT JOIN
         3 TABLE ACCESS FULL ARTISTE
       4 SORT JOIN
         5 TABLE ACCESS FULL FILM
\end{minted}

\end{frame}

\begin{frame}{Plan d'exécution}

\figSlideWithSize{planOra-joinsansindex}{8cm}

\end{frame}